\chapter{Conclusions \label{ch7}}

This doctoral research explored the use of Augmented Reality technologies to create a learning environment based on the principes of Collaborative Learning. \acf{AR}
technologies provide unique characteristics that can be used to promote diverse learning scenarios in which students participate more, report more engagement and develop
different dynamics with the learning material. Additionally, \acf{CL} approaches use social tools to encourage the students to gradually build the knowledge they need as
part of a community and better mirroring the environments and challenges that can be found in real life.

Of the characteristics identified, this project focused on exploring these main attributes that distinguish \ac{AR} from other platforms and \ac{CL} from other learning
philosophies:

\begin{itemize}
    \item Digital visuals that incorporate and react to the physical context of the user.
    \item Interactions with the environment that promote creative approaches to the learning material, exploration and active participation.
    \item Multiple communication channels that promote constant dialogue and conversation between participants.
    \item Guidance in identifying and using the skills necessary to create a collaborative effort as part of a group.
\end{itemize}

This project proposed that combining different characteristics of \ac{AR} with the connectivity and social tools provided by mobile platforms, it is possible to create a
learning space that promoted and facilitated the use collaboration and teamwork as teaching tools and learning goals. The main objective was to identify how students
approached challenges of collaborative learning using digital technologies, and to design a tool based on \ac{AR} that could promote behaviours that boost or reinforce
effective collaboration and the development of teamwork skills.

The approach of the project to fulfil this objective was divided into two parts. First, a software developed using user-centred design that identified core issues about
interactions in \ac{AR} spaces, multiuser experiences using \ac{AR}, and the challenges associated to develop software for learning support. The second part was the
analysis of a case study in which \ac{AR} was used as a tool to provide support for the objectives of a course that were using a collaborative learning philosophy or
developing collaboration as a skill.

\emph{WorkshopAR} was the prototype software born from these two research parts and was used to test different features that used \ac{AR} interactions to help students manage a
collaborative activity, while supporting the course in teaching basic principles of effective teambuilding and cooperation practices. \emph{WorkshopAR} provided tools for:

\begin{itemize}
    \item Augmenting the space that a group of students is using to work, providing visuals and interactions to help with team consolidation, team management, planification,
    progress evaluation, work management, generation of ideas, knowledge management, decision making and participation management.
    \item Creating a structured approach to teamwork that can be configured to fit the different stages of a project and the different characteristics and goals a group can have.
    \item Augmented communication for face-to-face interactions, using visuals and interactions to communicate ideas, expand the context of a conversation, provide extra
    meaning or to explore ideas and be creative.
\end{itemize}

The design process of \emph{WorkshopAR} identified important challenges in the implementation of \ac{AR} software and the development of technologies for learning settings.
It highlighted the importance of early testing and of adapting a flexible development approach to identify the particularities of a classroom, understand the needs of the
target students and adapt to the different profiles that could be found among them. A classroom is never homogeneous and is never static, and the implementation of the
software should be able to adapt to these variable characteristics.

The Industry Project course provided an ideal case study to observe a collaborative design. The core objective was to offer the opportunity to experience a project close to
a real-life professional setting, encouraging students to develop and test the skills needed for teamwork, communication, planning and constant self-reflection.

\emph{WorkshopAR} was integrated into the course as part of the repository of resources available to the students to help them manage the face-to-face collaborative
activities present in the classroom sessions. The students found opportunities to use \emph{WorkshopAR} by:

\begin{itemize}
    \item Using the structure provided by the tool at convenience, analysing the current state of their project, the objectives set up for their work, and the disposition
    of the current participants.
    \item Using the space and resources available in the current context to propose ideas, promote discussions and take decisions.
    \item Maintaining constant communication with their teammates, not only transmitting information but also contextual meaning.
\end{itemize}

The development of the case study showed the importance of giving the students the space to create collaborative skills in an organic way but also offering guidance when
needed, mechanisms to fail safely and opportunities to iterate in the approaches proposed by a group. \emph{WorkshopAR} also benefited from this iterative approach,
failing to integrate in some aspects of the collaborative process but adapting, incorporating feedback and finally finding an interesting place in other instances of the
students' work.

\section{Research Questions \label{sec:7-1}}

The research structure of this doctoral project focused on identifying changes in behaviour produced by the \ac{AR} tool supporting the collaborative work of the students.
\ac{AR} was used in activities related to team management, communication, and decision making.

This thesis proposed six secondary research questions to help build a conclusion for the main objective of the project. The questions were aimed at comparing current
approaches and perceptions that students have around collaboration and group-based learning against what \ac{AR} could offer. The main objective of the software
development was to test ideas in which the technology could create a different approach than those already used by the students.

\begin{displayquote}
    SQ1: What differences in behaviour can be observed between students collaborating using the \ac{AR} tool and those that are not?
\end{displayquote}

The use of \emph{WorkshopAR} promoted the most positive behaviours at the beginning of the life cycle of the project, when the students were first consolidating the group,
identifying the collaborative mechanisms to be used in the activity and deciding what rules and what roles would be present in the dynamics of their group.

Students were discussing how to organise their group while configuring the \emph{digital room} and identifying elements like the session goal and the assignation of tasks.
Even when not directly using \emph{WorkshopAR}, discussions over these elements continued after their introduction, and helped the students with the organisation of their
work.

\emph{WorkshopAR} also created a positive influence with the quick interactions taking place between students during the work session like conversations, planning and
decision making. The bulk of the communication of the group was taking place outside of the classroom, and major instances of communication like team organisation and the
flow of information needed for the project was mediated by other technologies and taking place using asynchronous methods. \emph{WorkshopAR} gave the students tools to have more
dynamic conversations during the classroom time, which was observed to be primarily focused on individual work. The students were using visuals and low structured debates
to provide more information about their activities, to add more context to the conversations taking place or simply interacting more between them, which was closer to the
initial intent that the teaching staff had for those presential sessions.

The most relevant negative behaviour observed promoted by \emph{WorkshopAR} was an excessive allocation of time in configuration tasks. The work of flow of the groups was
interrupted every time that someone wanted to interact with some feature of \emph{WorkshopAR}. If a long time had passed between interactions, the configuration or setup
needed was even longer. The work of the students was characterised by medium to long periods of intensive work using a variety of technologies and devices, and \emph{WorkshopAR}
failed to integrate well with this environment, disincentivising its use in favour of quicker or more effective means of communication.

\begin{displayquote}
    SQ2: What differences do the students report in their perception of the activities done in the course between those using the \ac{AR} tool and those who are not? 
\end{displayquote}

One of many requirements for the students to create an effective collaboration during the course was to identify and coordinate all the software tools needed to work and
communicate. Students saw \emph{WorkshopAR} as another tool in the repertoire of options at their disposal, and not in contraposition to any other. Some students saw the
value in the characteristics of the tool and incorporated it in their work, others preferred to use tools they felt more comfortable with.

Students were able to analyse if the use of \emph{WorkshopAR} was beneficial for a work session or not and setting aside the tool was as valid a decision as trying to use
the features of \emph{WorkshopAR}. Rather than changing perceptions or behaving differently with the use of one tool or another, some students incorporated WorkshopAR into
their workflow and others not, and both were correct approaches to the situation.

\begin{displayquote}
    SQ3: How does the collaboration and communication process changes between students using the \ac{AR} tool and those that are not? 
\end{displayquote}

\emph{WorkshopAR} helped in creating more dynamic and meaningful forms of conversation in the face-to-face setting of the course, as can be evidenced in the observations
detailed in Sections \ref{sec:6-5-2-1} and \ref{sec:6-8-2-2}.

The collaborative design followed by the Industry Project course presented the students with the challenge of understanding the problem of the project and proposing a
solution as a group, but relied mostly on individual work for the creation of the deliverables and the assessment of the project, incentivising minimal communication among
the members of the group once certain milestones of the project were reached.

The groups that used \emph{WorkshopAR} had more opportunities and resources to communicate between each other during the work session. They discussed more elements about
the structure of their work and the current progress of the project and were more open to discuss their individual work with their peers. Those conversations were short
but consistent, and their working session were more dynamic.

In contrast, students not using \emph{WorkshopAR} focused more on individual work and the use of personal devices. Communication was mostly with lecturers or with a \acf{TA},
and the information gathered benefited only their individual work. It did not mean that communication was not present, or that their team management was poor. It meant that
those instances of communication were taking place asynchronously and outside the classroom, and the face-to-face session were not being used to their full potential.

\begin{displayquote}
    SQ4: What differences can be observed in the end results of the course between all the different collaborative approaches observed, especially of those that incorporate the
    \ac{AR} tool? 
\end{displayquote}

There was no evidence of a relation between the use of \emph{WorkshopAR} and the results achieved by the students. Groups that presented good results were those that organised their
resources and established a proper workflow, which sometimes included the use of \emph{WorkshopAR}. Groups that had a better sense of the structure of their work used the tool more for
conversations and debates, while groups that initially struggled more organising themselves used the tool to as an aide and then moved  to other objectives. \emph{WorkshopAR} was present
in both successful and struggling projects.

The last two secondary questions changed focus to the implementation of the tool itself rather than its role in the case study and required the analysis of what elements of \ac{AR} were
creating the behaviours observed and how they could be better understood and incorporated in future projects.

\begin{displayquote}
    SQ5: What elements of a technological intervention using \ac{AR} create the most impact in changing the behaviours, perceptions or results of the students? 
\end{displayquote}

The \ac{AR} implementation used by \emph{WorkshopAR} succeeded in conveying information about structure and effective team management as well as engaging the students with dynamic and
constant communication when other groups were struggling to be more active during the classroom sessions.

In contrast, \emph{WorkshopAR} failed to offer value to the students during activities related to the design and implementation of the proposal for their project. The tool required long
configurations and the interactions offered through the play models were only useful during the initial stages of the project.

\emph{WorkshopAR} was good for proposing and discussing ideas but not to support co-creative collaboration or help in the work needed for the deliverables of the project. The augmentation
core of \emph{WorkshopAR} worked better when conveying information in a visual format to the students, but the physical interactions with the digital space were clunky and consumed too
much time. In either case, the objective to provide an augmented layer of information to the physical space and to the communication process can be considered positive as shown in Sections
\ref{sec:5-2-3} and \ref{sec:6-5-2}. Interesting points of discussion for future projects can also be extracted from the analysis of SQ5:

\begin{itemize}
    \item Data annotation linked both to the physical space and to the \emph{digital room} where participants are interacting. Offer the possibility to save this information, modify it
    across multiple sessions, and share it with other similar digital spaces.
    \item More varied methods to support face-to-face communication using visuals and shared interactions, both in structured and unstructured settings.
    \item Connectivity and synchronisation, not only with the tools provided by mobile platforms (such as photos, contact list and web browsing), but with other tools that are often used
    in collaboration settings like group chats, repositories of ideas and sketch tools.
\end{itemize}

\begin{displayquote}
    SQ6: What elements of a multi-user \ac{AR} tool have the most impact in changing the behaviours, perceptions, or results of students collaborating and communicating in a
    learning scenario? 
\end{displayquote}

The possibility of sharing a digital space with the group was at the core of the features proposed by \emph{WorkshopAR}, and it resonated well with the students. The shared space offered
tools to represent and give a more tangible presence to the group and allowed the students to express their personality using the tools present in the User Anchors and the emotes. This type
of social activity helped in the dynamics of the groups and promoted a better use of the presential work sessions.

Students were also organised around asynchronous forms of communication and work that WorkshopAR was not supporting, and the students found very limiting that there was no way to coordinate
those channels, as both were crucial for their work.

On a similar note, \emph{WorkshopAR} did not incorporate the horizontal interactions between groups that the shared space on the classroom was promoting, nor did it integrate the role of
the teaching staff in its workflow, and it was evident that the interaction with WorkshopAR felt disjointed from the rest of the activities because of that.

The secondary research questions offered an overview of what was achieved with the development of WorkshopAR and its implementation in the Industry Project course. \ac{AR} was used to create
a shared environment supporting the activities taking place in the classroom. Observations in Sections \ref{sec:6-4-2} and \ref{sec:6-8-2-2} show how the technology can help students to
better develop the skills necessary for teamwork and collaboration. The participating groups used this tool the most to help them to consolidate their group and to engage in more dynamic
processes of conversation. With this information at hand, we can now revisit the main question of this research:

\begin{displayquote}
    MQ: How does the behaviour of a group of students changes when using an AR tool designed to aid in the achievement of a learning goal framed in a collaborative learning scenario?
\end{displayquote}

Many of the groups participating in the course showed a bad or poor approach  to the face-to-face collaboration setting proposed for the classroom sessions, overfocusing on individual
assignments and avoiding almost any form of communication within the group. \emph{WorkshopAR} offered the tools to create more dynamic spaces for conversation and integrated those features
better to the workflow proposed by the students, rather than forcing the students to follow a specific structure.

The Industry Project course provides different ways to help the students further develop collaborative skills, but mostly opted to let them explore by themselves the necessities of their
group and to find solutions to the hurdles they were encountering. It was satisfactory to see that \emph{WorkshopAR} helped some groups to identify what was lacking in their approach and to
continue the conversation of team management using the concepts and mechanisms that \emph{WorkshopAR} provided.

Students were also capable of analysing the resources at their disposals, especially those digital tools needed to complete their work. Students were able to discuss and recognise when
\emph{WorkshopAR} was not a good match for their work style or the activities planned for a session and took the decision to explore alternatives that better suited them as a group.

Students disliked the long configuration process required for many of the features of \emph{WorkshopAR} and tried to explore alternatives to either create a more complex working station
using the smartphone as a basis, or dedicating more time to the activities mediated by \emph{WorkshopAR} to justify the time invested in configuration. Neither of these approaches related
well with the original design of the tool or with the workflow proposed by the course and was affecting the time allocation that students had to manage during the classroom sessions in a
negative way.

In conclusion, it was possible to observe both positive and negative behaviour changes in the students that interacted with \emph{WorkshopAR}. The tool provided a different channel of
communication in an activity that the students were struggling to implement and offered learning tools to improve the collaborative skills during the early stages of group consolidation.
The tool was also consuming too much time in configurations and that was affecting negatively the time management process of the group.

Overall, it was shown that the unique characteristics of \ac{AR} can be used to implement an effective collaborative learning approach that promotes interesting forms of communication and
learning, but also poses substantial challenges based on the hardware used to deploy the solution and how much the technology has to integrate with the academical work of the students. Both
positive and negative results are valuable lessons that can further the research in education or to be a starting point to explore other forms of implementation of this technology.

\section{Future Work \label{sec:7-2}}

\emph{WorkshopAR} was developed following an iterative process aimed at polishing features in each cycle and to explore the ideas found by gathering feedback from the users. Through this
procedure, adjustments to the inconsistencies in the interaction design were identified and implemented, and the emotes feature was expanded when it became a focus of interest for the students.

Additional iterations over the development of \emph{WorkshopAR} would continue with this process of expansion of the features that students found compelling but were not complete enough.
Specifically, it could be of interest to provide the following expansions:

\begin{itemize}
    \item Annotations using other forms of multimedia beside images and text, potentially expanding the configuration process to incorporate content that could be found on the local storage
    of the device or on the web.
    \item Reimagine the play models feature to become a tool for augmented presentations. This idea came from conversations taking place with the students when asked how to make the play models
    tool more relevant for their work in the course (see Sections \ref{sec:6-5-2-2} and \ref{sec:6-6-2}). The interactions proposed included a shared presentation of three-dimensional information
    to explain the proposal of the group to the clients or the use of presentation slides and posters that could be projected and interacted with in the room. All these ideas hinted at a tool
    that supports a more structured form of communication and exposition of information.
    \item Incorporate the role of the teacher as initially proposed in the interaction design, offering tools to visualise and interact with the work of many groups at the same time and
    integrating the teacher and other similar roles into the workflow proposed by the students.
\end{itemize}

All these proposals are natural extension to the original design of \emph{WorkshopAR} that can be considered to create a more robust software or, following a different approach, could be individual,
smaller developments that branch from \emph{WorkshopAR}.

Other ideas and paths of research emerged from conversations with all the participants, especially when discussing how to solve the most evident issues with \emph{WorkshopAR} related to
long configurations and the constant use of the phone in the classroom.

The idea of a collaboration laboratory using \ac{AR} and other related technologies emerged from discussions with the students about how to create a setup that deemphasises the distractions
created by a smartphone and that promotes a better work mentality. The use of clamp arms in the setup of \emph{WorkshopAR} tried to create a space more suited for consistent collaboration
(see Sections \ref{sec:5-2-4-1} and \ref{sec:6-8-2-2}), and students expanded that idea to incorporate other forms of hardware and activities not necessarily linked with learning. This idea
tackles the problem found with excessive configuration requirements, but it abandons the premise of a \acf{BYOD} activity and the flexibility of augmenting any space for work. It could be
interesting to dig deeper into these ideas of fixed \ac{AR} spaces for collaboration that were difficult to test due to the characteristics of the case study.

A second scenario for development appeared when discussing how to expand the annotation feature. The idea took form when some students compared the way the annotations behaved to other web
tools like Miro, Milanote and Pinterest, and how students and professional use those tools to aid their creative process  (see Section \ref{sec:6-5-2-1}). This scenario would explore some
of the recommendations made by students for the annotations to expand their use into scenarios like visual repositories of ideas, mental maps, mood boards and other similar activities.

These seed ideas for future projects and for possible paths forward in the evolution of \emph{WorkshopAR} derived from conversations with the students and teaching staff, those that managed
to relate the functionalities of the software with the needs of their work, showing that \emph{WorkshopAR} was creating interest in the intended users. A progressive expansion and adjustment
of the proposal can result in a tool that is useful to the students and helps to promote a valuable learning design focused on teamwork and collaboration.

\section{Final Words \label{sec:7-3}}

The development of \emph{WorkshopAR} and its integration to the Industry Project course was a valuable example of a collaborative design and how this philosophy supports the creation of
transversal skills and complex learning scenarios by encouraging social connections, collective intelligence and effective communication.

It can be confidently said that this research did produce observable changes in the behaviours of the students incorporating the tool to the workflow that the course was encouraging. Students
demonstrated a better communication process and had the opportunities to further develop collaborative skills, and in turn they helped to evolve and grow the project itself thanks to the
feedback provided and the rich experiences that were possible to observe in the classroom, opening the possibilities to keep expanding this line of research in future work thanks to the
promising results obtained.